% =====================================================================
%  Section 2: Related Work
% =====================================================================
\section{Related Work}
\label{sec:related}

Instruction-tuning data selection has been studied through both
static and dynamic criteria. Quality-based methods assign each
candidate a global utility score, either using a strong external
judge, as in AlpaGasus~\cite{alpagasus}, or a learned quality
estimator, as in CPQS~\cite{cpqs2026}. These methods are simple and
task-agnostic, but a single pool-level ranking cannot adapt to a
target capability or to the model's changing state during training.

Target-conditioned methods define usefulness relative to examples
that represent a desired capability. Influence- and gradient-based
approaches such as LESS~\cite{less} score candidates by their
gradient similarity to a small validation or few-shot target set,
while Approx-LESS~\cite{approxless} improves scalability by
predicting influence scores from a subset of gradient features.
Another line uses representation geometry rather than gradients:
NAIT~\cite{nait2026} extracts hidden-state capability directions
from in-domain seed examples and selects candidates aligned with
those directions. These approaches can produce capability-aligned
subsets, but their selection criterion is fixed before training and
depends on seed or validation examples supplied by the practitioner.

\paragraph{Composite-reward dynamic selection.}
A separate line treats data selection as a sequential decision
problem and re-estimates per-candidate utility during training. The
composite-reward family~\cite{kung2023activeit, sun2025dots,
jha2025rl} combines loss-based difficulty $\mathcal{L}_i$ and
entropy-based uncertainty $\mathcal{H}_i$ into a single per-sample
scalar $\mathcal{R}_i = w\mathcal{L}_i + (1-w)\mathcal{H}_i$ with a
variance-ratio weight $w$, and ranks candidates by this scalar.
Some of these methods additionally train a policy network
(typically a PPO actor-critic) on top of the composite reward;
others use it directly. \textsc{TADS} adopts the composite-reward
construction unchanged as its base utility signal but applies a
deterministic pool-level normalisation and adds a representation-
level structural factor.

\paragraph{Dynamic data selection for RL fine-tuning.}
A closely related but experimentally non-comparable line studies
dynamic data selection for RL fine-tuning with verifiable rewards
(RLVR). DOTS~\cite{sun2025dots} targets GRPO training by selecting
questions whose \emph{adaptive difficulty} -- the policy-dependent
rollout failure rate -- is close to one half, where binary reward
variation yields the strongest group-relative gradient signal. DOTS
is therefore an effective dynamic question-selection method for
RLVR, but its signal is defined through rollout groups, verifiable
rewards, and a learned difficulty predictor.

DOTS and \textsc{TADS} share the principle of \emph{per-epoch
re-estimation of per-candidate utility}, but they are not
interchangeable baselines. DOTS selects RLVR questions under GRPO;
\textsc{TADS} re-ranks supervised instruction--response pairs under
SFT. Applying DOTS to our setting would require converting SFT data
into verifiable RLVR questions and adding reward design, rollout
generation, and difficulty prediction, thereby changing the
objective, data unit, supervision source, and compute profile. We
therefore treat DOTS as the closest dynamic selector in a different
training regime, not as a directly runnable baseline for dynamic
SFT data selection.

\paragraph{Positioning of \textsc{TADS}.}
The central contribution of \textsc{TADS} is the \emph{trajectory
anchor}: a per-layer capability direction extracted online from
the model's own hidden-state trajectory at each refresh step. The
anchor is what makes \textsc{TADS} representation-aware where
prior composite-reward dynamic selectors are not, and what
distinguishes it from static representation-based methods whose
capability directions are fixed before training and depend on seed
or target examples. Unlike static directions, the anchor is
re-extracted from the candidate pool itself at every epoch as the
model's hidden states evolve; unlike scalar composite rewards, it
captures direction in representation space rather than only
magnitude.

Around this anchor, \textsc{TADS} reuses the standard
composite-reward signal $\mathcal{R}_i^{(t)} = w^{(t)}\mathcal{L}_i^{(t)} + (1-w^{(t)})\mathcal{H}_i^{(t)}$ from prior
dynamic-selection work~\cite{kung2023activeit, sun2025dots} in a
supporting role, without any extra transform. The selection score
is the unmodified composite reward multiplied by the
trajectory-anchor factor:
\[
s_i^{(t)}
=
\mathcal{R}_i^{(t)}
\left(
1+\lambda\,\widetilde{\mathrm{align}}_i^{(t)}
\right),
\]
in which the trajectory-anchor factor on the right is the novel
structural component; the composite-reward factor on the left
is the per-candidate magnitude signal carried over from prior
work. Setting $\lambda=0$ disables the anchor and recovers the
composite-reward ranking alone, yielding a clean ablation that
isolates the anchor's contribution. Crucially, the trajectory
anchor is defined independently of $\mathcal{R}_i^{(t)}$; the
stability guarantee we prove in Sec.~\ref{sec:method-theory}
therefore concerns the structural component introduced by
\textsc{TADS}, and applies to any uniformly bounded base ranking
score on the left.
